% Einheit 3 von 4 – Rechtwinklige Teildreiecke in Figuren und Vermessung (bank/trigonometrie/e3.jsonl, zusammenbau v0.1)
%% TODO Zweigzeile Teil 1: was hier gelernt wird (Satz in Schülersprache aus dem Einheitstitel) – Einheit 3
\einheitenkopf[e3]{Einheit 3 von 4 · Rechtwinklige Teildreiecke in Figuren und Vermessung}
\zweigzeile{ab Kl. 10 (am Gymnasium ab Kl. 9) · P10 oft}
\setcounter{aufgabe}{22}

%% TODO Titel: Ich-kann-Satz fehlt in der Bank (Platzhalter: Kettenname) – Nr. 23
%% TODO Anweisung: ein Satz über der ersten Teilaufgabe fehlt in der Bank – Nr. 23
\begin{aufgabe}{Teildreiecke und Vermessung}
%% TODO \rechenplatz{n}: Schrittzahl des Grundfalls fehlt in der Bank (nur bei fester Schreibform, 2.3 b)
\begin{teile}
\teil Fahre im Trapez das rechtwinklige Teildreieck mit der Höhe $h$ links nach. Beschrifte seine Seiten mit H, G, A – vom Basiswinkel unten links aus.

\trapez[hoehe=h]{5}{3}{2.5}
\teil Gleichschenkliges Dreieck: Schenkel $7$ cm, Basiswinkel $63^\circ$. Wie lang ist die Höhe $h$ auf die Grundseite? h ≈ \leerfeld[cm]
\teil Trapez: Höhe $4$ m, Basiswinkel $62^\circ$. Wie lang ist der Schenkel $s$? s ≈ \leerfeld[m]
\teil Trapez: Schenkel $5$ cm, Basiswinkel $58^\circ$. Wie hoch ist das Trapez? h ≈ \leerfeld[cm]
\teil Rechtwinkliges Trapez: die senkrechten Seiten sind $9$ cm und $4$ cm lang und stehen $8$ cm auseinander. Wie groß ist der Winkel $\alpha$ zwischen der schrägen Seite und der Waagerechten? $\alpha \approx$ \leerfeld $^\circ$
\teil Parallelogramm $ABCD$: $BC = 9$ cm. $F$ liegt auf der Verlängerung von $AB$ über $B$ hinaus, $CF$ steht senkrecht darauf; der Winkel $CBF$ beträgt $58^\circ$. Wie lang ist die Höhe $h = CF$? \leerfeld[cm]
\teil Parallelogramm $ABCD$ mit der Höhe $h = 6$ cm von $C$ auf die Verlängerung von $AB$. Die Diagonale $AC$ bildet mit $AB$ einen Winkel von $27^\circ$. Wie lang ist $AC$? AC ≈ \leerfeld[cm]
\teil Drachen: eine obere Seite ist $6$ cm lang und bildet mit der Längsdiagonale einen Winkel von $32^\circ$. Wie lang ist die halbe Querdiagonale? \leerfeld[cm]
\teil Aus $45$ m waagerechtem Abstand sieht man die Spitze eines Baums unter einem Höhenwinkel von $31^\circ$. Wie hoch liegt sie über dem Messgerät? \leerfeld[m]
\teil Ein Messgerät steht in $1{,}4$ m Höhe, $58$ m waagerecht entfernt. Der Höhenwinkel zur Spitze eines Kirchturms beträgt $27^\circ$. Wie hoch ist es? \leerfeld[m]
\teil Aus $48$ m Abstand peilt man in $1{,}7$ m Augenhöhe die Dachkante eines Turms unter $24^\circ$ an. Auf dem Dach steht eine $6$ m hohe Antenne. Wie hoch liegt die Antennenspitze über dem Boden? \leerfeld[m]
\teil Rechter Winkel bei $D$; $B$ liegt zwischen $D$ und $C$. Von $A$ aus: $\angle DAC = 71^\circ$, $\angle BAC = 23^\circ$, $AD = 32$ m. Wie lang ist $BD$? BD ≈ \leerfeld[m]
\teil Dreieck $ABC$ mit der Höhe $h$ von $C$ auf $AB$: $AC = 8$ cm, $\alpha = 42^\circ$, $\beta = 61^\circ$. Berechne erst $h$, dann $BC$. h ≈ \leerfeld[cm] , BC ≈ \leerfeld[cm]
\teil Gleichschenkliges Trapez: bekannt sind die beiden parallelen Seiten und die Höhe, gesucht ist der Umfang. Beschreibe den Rechenweg ohne zu rechnen.
\teil Hypotenuse $9$ cm, $\alpha = 33^\circ$. Berechne Gegenkathete und Ankathete und prüfe mit dem Satz des Pythagoras. G ≈ \leerfeld[cm] , A ≈ \leerfeld[cm]
\teil Dreieck $ABC$: $c = 9$ cm, $b = 6$ cm, $\alpha = 47^\circ$. Berechne die Höhe $h$ auf $c$ und dann den Flächeninhalt. h ≈ \leerfeld[cm] , A ≈ \leerfeld[cm²]
\teil Quadratische Pyramide: Grundkante $6$ m, die Seitenflächen sind um $57^\circ$ gegen die Grundfläche geneigt. Wie lang ist die Seitenhöhe $h_s$? \leerfeld[m]
\teil Ein Messgerät steht $1{,}4$ m über dem Boden. Die Sichtlinie zur Hügelkuppe ist $96{,}5$ m lang und steigt unter $27^\circ$ an. Wie hoch ist der Hügel? Auf der Kuppe steht ein $9$ m hoher Aussichtsplattform – wie hoch über dem Boden liegt seine Oberkante? \hfill (P10 2018 OS) \\ \leerfeld[m] ; \leerfeld[m]
\end{teile}
\end{aufgabe}

%% TODO Titel: Ich-kann-Satz fehlt in der Bank (Platzhalter: Kettenname) – Nr. 24
%% TODO Anweisung: ein Satz über der ersten Teilaufgabe fehlt in der Bank – Nr. 24
\begin{aufgabe}{Strecke aus Teilstrecken nach der Trigonometrie}
\begin{teile}
\teil Dreieck $ABC$ mit der Höhe $CD = 4$ cm auf $AB$: $AD = 6$ cm, $\beta = 53^\circ$. Wie lang ist $AB$? AB ≈ \leerfeld[cm]
\end{teile}
\end{aufgabe}

%% TODO Titel: Ich-kann-Satz fehlt in der Bank (Platzhalter: Kettenname) – Nr. 25
%% TODO Anweisung: ein Satz über der ersten Teilaufgabe fehlt in der Bank – Nr. 25
\begin{aufgabe}{Teildreiecke und Vermessung – Fehler finden, Begründen, Anwendung}
\begin{teile}
\teil Abstand zum Turm $44$ m, Höhenwinkel $31^\circ$, das Messgerät steht $1{,}7$ m hoch. Wo ist der Fehler? \rechnung{h &= 44 \cdot \mathrm{tan}\,31^\circ \\ h &\approx 26{,}4 \text{ m}}
\teil Warum rechnet man im gleichschenkligen Dreieck mit der halben Grundseite, wenn man den Basiswinkel benutzt?
\teil Ein Baum soll gefällt werden, das Haus steht $18$ m vom Stamm entfernt. Aus $20$ m Abstand sieht man die Baumspitze in $1{,}7$ m Augenhöhe unter $41^\circ$. Kann der Baum beim Fallen das Haus treffen?
\end{teile}
\end{aufgabe}
